\documentclass[11pt,a4paper]{article}

%% PACHETE MATEMATICE
\usepackage{amsmath,amssymb,amsfonts}
\usepackage{amsthm}
\usepackage{mathpazo}
\usepackage{eulervm}
\usepackage{enumerate}
\usepackage{color}
\usepackage{graphicx}
\usepackage[all]{xy}
\usepackage{xcolor}
\usepackage{framed}
\usepackage[unicode]{hyperref}
\setcounter{MaxMatrixCols}{20}
\usepackage{multicol}
%\usepackage[romanian]{babel}


%% ASPECT
\linespread{1.05}
\usepackage[T1]{fontenc}
\usepackage[utf8x]{inputenc}
\usepackage[margin=2cm]{geometry}
% \usepackage[color]{showkeys}
\usepackage{anyfontsize}
\usepackage{titlesec}
\titleformat*{\section}{\large\bfseries}

\usepackage{fancyhdr}
\pagestyle{fancy}
\rhead{\small \color{gray} Notițe de Adrian Manea}
\lhead{\small \color{gray} IntProb-FIA, 2017---2018, L. Lemnete \& M. Olteanu}


\usepackage[autostyle = false, style = german]{csquotes}
\usepackage[romanian]{babel}
\newcommand\qq{\enquote}

%% COMENZILE MELE
\renewcommand*{\proofname}{Dem.:}
\newcommand\bb{\mathbb}
\newcommand\dr{\mathrm}
\newcommand\kal{\mathcal}
\newcommand\fk{\mathfrak}
\newcommand\op{\oplus}
\newcommand\ot{\otimes}
\newcommand\ds{\displaystyle}
\newcommand\id{\indent}
\newcommand\nid{\noindent}
\newcommand\seq{\subseteq}
\newcommand\sneq{\subsetneq}
\newcommand\speq{\supseteq}
\newcommand\spneq{\supsetneq}
\newcommand\rar{\rightarrow}
\newcommand\Rar{\Rightarrow}
\newcommand\xrar{\xrightarrow}
\newcommand\lar{\leftarrow}
\newcommand\Lar{\Leftarrow}
\newcommand\xlar{\xleftarrow}
\newcommand\lrar{\leftrightarrow}
\newcommand\Lrar{\Leftrightarrow}
\newcommand\ideal{\trianglelefteq}
\newcommand\ai{\text{ a.\^{i}. }}
\newcommand\sk{(\textit{Schi\c{t}\u{a}}) }
\newcommand\rhup{\rightharpoonup}
\newcommand\rhdn{\rightharpoondown}
\newcommand\lhup{\leftharpoonup}
\newcommand\lhdn{\leftharpoondown}
\newcommand\vid{\emptyset}
\newcommand\wh{\widehat}
\newcommand\ol{\overline}
\newcommand\NN{\mathbb{N}}
\newcommand\RR{\mathbb{R}}
\newcommand\ZZ{\mathbb{Z}}
\newcommand\QQ{\mathbb{Q}}
\newcommand\CC{\mathbb{C}}

%% STILURILE MELE
\newtheoremstyle{dotless}{}{}{\itshape}{}{\bfseries}{:}{ }{}

\theoremstyle{dotless}
\newtheorem{theorem}{Teorem\u{a}}[section]
\newtheorem{proposition}{Propozi\c{t}ie}[section]
\newtheorem{lemma}{Lem\u{a}}[section]
\newtheorem{corollary}{Corolar}[section]
\newtheorem{exercise}{Exerci\c{t}iu}

\newtheoremstyle{dotlessdr}{}{}{}{}{\bfseries}{:}{ }{}
\theoremstyle{dotlessdr}
\newtheorem{example}{Exemplu}[section]
\newtheorem{definition}{Defini\c{t}ie}[section]
\newtheorem{remark}{Observa\c{t}ie}[section]




\begin{document}

\begin{center}
  {\large\textbf{Seminar 1 \\ Spații metrice}}
\end{center}

\vspace{2cm}


\section{Spații metrice}
\label{sec:sp-metrice}

\begin{definition}\label{def:distanta}
  Fie $X$ o mulțime nevidă. O aplicație $d : X \times X \to \RR$ se numește \emph{distanță}
  (metrică) pe $X$ dacă:
  \begin{enumerate}[(a)]
  \item $d(x, y) \geq 0, \forall x, y \in \RR$;
  \item $d(x, y) = 0 \Leftrightarrow x = y$;
  \item $d(x, y) = d(y, x), \forall x, y \in X$;
  \item $d(x, z) \leq d(x, y) + d(y, z)$ (\emph{inegalitatea triunghiului}).
  \end{enumerate}

  În acest context, perechea $(X, d)$ se numește \emph{spațiu metric}.
\end{definition}

Această noțiune vine să generalizeze calculul distanțelor cu ajutorul modulului,
cum se procedează în cazul mulțimii numerelor reale, de exemplu. În consecință,
avem și următoarele \emph{mulțimi remarcabile în spații metrice $(X, d)$}:
\begin{itemize}
\item \emph{bila deschisă de centru $a$ și rază $r$}, definită pentru $a, r \in \RR$ prin:
  \[
    B(a, r) = \{ x \in X \mid d(a, x) < r \};
  \]
\item \emph{sfera de centru $a$ și rază $r$}, definită prin:
  \[
    S(a, r) = \{ x \in X \mid d(a, x) = r \};
  \]
\item \emph{bila închisă de centru $a$ și rază $r$}, definită prin:
  \[
    \overline{B(a, r)} = B(a, r) \cup S(a, r) = \{x \in X \mid d(a, x) \leq r \}.
  \]
\end{itemize}

Aceste mulțimi generalizează \emph{intervalele de numere reale}. Mai general, avem:
\begin{definition}\label{def:m-deschisa}
  Fie $X$ un spațiu metric. O submulțime $D \seq X$ se numește \emph{deschisă} dacă
  $\forall a \in D, \exists r > 0$ astfel încît $B(a, r) \seq D$. Cu alte cuvinte,
  mulțimea conține o bilă deschisă, centrată în orice punct al său.

  O mulțime se numește \emph{închisă} dacă complementara ei, relativ la spațiul total,
  este o mulțime deschisă.
\end{definition}

Adaptînd noțiunile specifice pentru analiza matematică, precum convergență,
limită ș.a.m.d. folosindu-ne de funcția distanță, putem defini construcțiile
și conceptele uzuale. În particular, au sens noțiuni precum \emph{șiruri convergente}
și \emph{șiruri Cauchy}, definite ca în cazul mulțimii numerelor reale, doar
cu ajutorul funcției generale distanță. În plus, avem și următoarea noțiune:

\begin{definition}\label{def:spatiu-complet}
  Un spațiu metric $(X, d)$ se numește \emph{complet} dacă noțiunile de
  \qq{șir Cauchy} și \qq{șir convergent} coincid pentru $X$.
\end{definition}

%%%%%%%%%%%%%%%%%%%%%%%%%%%%%%%%%%%%%%%%%%%%%%%%%%%%%%%%%%%%%%%%%%%%%%

\section{Principiul contracției. Metoda aproximațiilor succesive}
\label{sec:contractie}

\begin{definition}\label{def:contractie}
  Fie $(X, d)$ un spațiu metric și fie $f : X \to X$ o funcție.

  Aplicația $f$ se numește \emph{contracție pe $X$} dacă există $k \in [0, 1)$
  astfel încît:
  \[
    d(f(x), f(y)) \leq k \cdot d(x, y), \forall x, y \in X.
  \]

  Numărul $k$ se numește \emph{factor de contracție}.
\end{definition}

\begin{remark}\label{rk:contr-der}
  De remarcat că în cazul clasic, unde $d(x, y) = |x - y|$, formula contracției
  se poate gîndi ca o aproximație finită a derivatei. Mai precis, avem:
  \[
    \frac{d(f(x), f(y))}{d(x, y)} = \frac{|f(x) - f(y)|}{|x - y|} \leq k.
  \]

  Așadar, calculînd $\sup |f'(x)|$ putem găsi valoarea maximă a factorului
  de contracție.
\end{remark}

Un rezultat important este următorul:

\begin{theorem}[Banach]\label{thm:banach-fix}
  Fie $(X, d)$ un spațiu metric complet și fie $f : X \to X$ o contracție de factor
  $k$. Atunci există un unic punct $\xi \in X$, astfel încît $f(\xi) = \xi$.

  În acest context, $\xi$ se numește \emph{punct fix} pentru $f$.
\end{theorem}

Pentru a găsi punctul fix al unei aplicații, se folosește \emph{metoda %
  aproximațiilor succesive}. Se construiește un șir recurent astfel.

Fie $x_0 \in X$, arbitrar și definim șirul recurent $x_{n+1} = f(x_n)$.
Se poate demonstra că șirul $x_n$ este convergent, iar limita sa este punctul
fix căutat. În plus, eroarea aproximației cu acest șir este dată de:
\[
  d(x_n, \xi) \leq \frac{k^n}{1 - k} \cdot d(x_0, x_1), \forall n \in \NN.
\]

%%%%%%%%%%%%%%%%%%%%%%%%%%%%%%%%%%%%%%%%%%%%%%%%%%%%%%%%%%%%%%%%%%%%%%

\section{Exemplu rezolvat}
\label{sec:ex-rez}

Aplicațiile interesante pentru această temă sînt date de calculul aproximativ
al soluțiilor unor ecuații, definite în spații metrice.

1. Să se aproximeze cu o eroare mai mică decît $10^{-3}$ soluția reală a
ecuației $x^3 + 4x - 1 = 0$.

\emph{Soluție:} Folosind, eventual, metode de analiză (șirul lui Rolle),
se poate arăta că ecuația are o singură soluție reală $\xi \in (0, 1)$.
Folosim metoda aproximațiilor succesive pentru a o găsi.

Fie $X = [0, 1]$ și $f: X \to X, f(x) = \frac{1}{x^2 + 4}$. Ecuația dată este
echivalentă cu $f(x) = x$, adică a găsi un punct fix pentru funcția $f$.

Spațiul metric $X$ este complet. Mai demonstrăm că $f$ este contracție
pe $X$. Derivata este $f'(x) = \frac{-2x}{(x^2 + 4)^2}$, și avem:
\[
  \sup_{x \in X} | f'(x) | = -f'(1) = \frac{2}{25} < 1,
\]
deci $f$ este o contracție, de factor $k = \frac{2}{25}$.

Șirul aproximațiilor succesive este dat de:
\[
  x_0 = 0, \quad x_{n + 1} = f(x_n) = \frac{1}{x_n^2 + 4}.
\]

Evaluarea erorii:
\[
  |x_n - \xi| < \frac{k^n}{1 - k}|x_0 - x_1| = \frac{1}{3} \Big( \frac{2}{25} \Big)^n,
\]
de unde rezultă că putem lua $\xi \simeq x_3 = f \big(\frac{16}{65}\big) \simeq 0,235$.

\begin{remark}\label{rk:alternativ}
  Alternativ, puteam lucra cu $g(x) = \frac{1}{4} (1 - x^3)$, cu $x \in [0, 1]$.
  Se arată că și $g$ este o contracție, de factor $k = \frac{3}{4}$. În acest caz,
  șirul aproximațiilor succesive converge mai încet și avem $\xi \simeq x_6$.
\end{remark}

\vspace{1cm}

Similar, puteți rezolva următoarele ecuații, cu eroarea $\varepsilon$:
\begin{enumerate}[(a)]
\item $x^3 + 12x - 1 = 0, \quad \varepsilon = 10^{-3}$;
\item $x^5 + x - 15 = 0, \quad \varepsilon = 10^{-3}$;
\item $3x + e^{-x} = 1, \quad \varepsilon = 10^{-3}$;
\item $x^3 - x + 5 = 0, \quad \varepsilon = 10^{-2}$;
\item $x^5 + 3x - 2 = 0, \quad \varepsilon = 10^{-3}$.
\end{enumerate}


\end{document}